\subsection{The Mahan Hypothesis: Decisive Battle and Sea Power}

Alfred Thayer Mahan's \textit{The Influence of Sea Power upon History, 1660--1783} \citep{mahan1890} established the theoretical foundation for the decisive battle hypothesis. Drawing on Britain's naval dominance in the eighteenth century, Mahan emphasized concentration of naval power, command of maritime communications, and the ability of sea power to influence national strategy as the primary determinants of war outcomes. The decisive battle, in Mahan's framework, is not merely a tactical victory but a strategic mechanism: by destroying the enemy's naval capability in a single engagement, the victorious power secures the conditions for political settlement on favorable terms. Critically, Mahan's analysis extended well beyond fleet engagements to encompass commercial power, maritime logistics, and the broader system of sea lines that sustained national power \citep{mahan1890}. Mahan recognized that naval dominance served economic and political objectives, not military ones in isolation.

Julian Corbett \citep{corbett1911} offered a complementary and in some respects more nuanced perspective, arguing that maritime strategy involves not only decisive fleet engagements but also limited objectives, sea control, and the projection of power ashore. Corbett emphasized the interaction between sea control and land operations, recognizing that most wars cannot be decided by naval action alone. His framework bridges the decisive battle and attrition traditions by showing how maritime power enables sustained operations while creating opportunities for decisive engagement when conditions permit. Later strategists such as Howard \citep{howard1979} and Luttwak \citep{luttwak1976} further enriched the theoretical landscape: Howard analyzed how Moltke's campaigns of annihilation through encirclement battles became the template for modern offensive warfare, while Luttwak revived interest in the decisive battle by arguing that the dialectical logic of strategy creates opportunities for decisive outcomes when an adversary makes an exploitable error.

The decisive battle tradition extends well beyond Mahan's naval focus. Clausewitz himself, despite being the intellectual forefather of attrition theory, recognized the importance of the \textit{Schlachtentscheidung}, the battle decision, as a means of breaking an enemy's will \citep{clausewitz1832}. Empirically, the decisive battle hypothesis has been applied to numerous historical cases. Wawro's analysis of the Franco-Prussian War \citep{wawro2003} emphasizes how the Battle of Sedan (1870) captured the French emperor and destroyed the main field army, effectively deciding the war's outcome even though fighting continued for several more months. The Gulf War of 1991 is frequently cited as a modern example, where a hundred-hour ground campaign decisively expelled Iraqi forces from Kuwait \citep{friedman1998}.

\subsection{The Attrition Hypothesis: Strategic Exhaustion}

The attrition hypothesis has roots in Clausewitz's observation that war is ``a continuation of politics by other means'' and that the destruction of an enemy's armed forces is merely a means to an end, not the end itself \citep{clausewitz1832}. In this view, wars are decided not by brilliant maneuvers but by the relative capacity of belligerents to sustain losses and continue fighting. The side with greater resources, deeper reserves, and stronger political will, or the side that can impose greater cumulative costs on its adversary, will ultimately prevail.

Modern attrition theory draws on several intellectual traditions. The ``culminating point'' concept, developed by Clausewitz and elaborated by Shy \citep{shy1976}, suggests that offensive operations reach a point of diminishing returns where the attacker's losses exceed the defender's, leading to exhaustion rather than decisive victory. Biddle's \textit{Military Power} \citep{biddle2004} provides a systematic framework for understanding how force employment (the way troops are actually used in combat) mediates the relationship between material resources and battlefield outcomes, often producing attritional dynamics even when commanders seek decisive results.

The quantitative study of attrition has been formalized through Lanchester's laws \citep{lanchester1916}, which model combat as a system of differential equations relating force sizes, attrition rates, and combat effectiveness. Subsequent work by Beckman \citep{beckman1995} and others has extended these models to incorporate stochastic elements, logistical constraints, and strategic decision-making. The ``loss exchange ratio'' (the ratio of casualties between opposing forces) has been extensively studied as a predictor of war outcomes, with consistent findings that wars with favorable exchange ratios tend to terminate on the superior side's terms \citep{heginbotham2015}.

\subsection{Previous Computational Approaches}

Computational approaches to studying war termination have proliferated in recent years. The Correlates of War (COW) project \citep{singer1972} provided the first systematic dataset of interstate conflicts, enabling quantitative analysis of war determinants. Subsequent projects, the UCDP/PRIO Armed Conflict Dataset \citep{gleditsch2002}, the Interstate War Battle Dataset \citep{rester2019}, and the Brecke Conflict Catalog \citep{brecke1999}, have expanded the temporal and geographic scope of conflict data.

Statistical analyses of these datasets have tested various hypotheses about war termination. Huth \citep{huth1996} examined the role of military balance in determining war outcomes, finding that relative power matters but that political factors mediate its effects. Bennett and Stam \citep{bennett2000} applied duration models to war length, identifying regime type and alliance structures as significant predictors. More recently, machine learning approaches have been applied to predict war outcomes from pre-war features, with mixed success \citep{ward2010}.

Despite these advances, the computational literature has largely treated war termination as a single outcome variable (victory, defeat, or negotiated settlement) rather than modeling the underlying dynamics that produce these outcomes. The decisive battle versus attrition question has been addressed indirectly through case studies and small-n comparisons \citep{reed2003}, but no large-n computational study has systematically quantified both shock and exhaustion dynamics across thousands of conflicts.

\subsection{Research Gap and Contribution}

Three gaps motivate the present study. First, existing metrics for classifying wars as ``decisive'' or ``attritional'' rely on subjective historical judgment rather than systematic quantitative criteria. Second, no study has developed complementary metrics for both decisive shocks and strategic exhaustion that can be applied consistently across a large and diverse dataset. Third, simulation models of war dynamics have not been systematically evaluated against historical data at scale.

This paper addresses these gaps by: (1) developing the DSS and SES metrics that quantify decisive shock and strategic exhaustion dynamics across thousands of wars; (2) applying these metrics to a comprehensive dataset of 4,812 conflicts; (3) constructing a dynamical simulation model that generates war trajectories under varying parameter regimes; and (4) evaluating the simulation against 7 historical case studies using a mechanism classifier that separates termination events from strategic causes. Our framework moves beyond the binary ``decisive or attritional?'' question to ask a more nuanced question: when and why do decisive shocks matter relative to accumulated exhaustion?
